\documentclass[uplatex,a4paper,9.5pt,twocolumn]{jsarticle}
\usepackage[margin=25.4mm]{geometry}
\usepackage{url}

% セクション見出しの上下余白のコンパクト化
\makeatletter
\renewcommand{\section}{\@startsection{section}{1}{\z@}%
  {0.4ex plus 0.1ex minus 0.1ex}%
  {0.2ex plus 0.05ex}%
  {\reset@font\normalsize\bfseries}}
\makeatother

% リスト環境の余白圧縮
\setlength{\leftmargini}{1.1em}
\setlength{\leftmarginii}{1.1em}

\setlength{\parindent}{1zw}
\setlength{\columnsep}{6mm}
\pagestyle{empty}

\begin{document}

% タイトル部分（コンパクト化）
\twocolumn[
  \begin{center}
    {\large\bfseries Spring BootとMySQLによるECサイト開発レポート}\par
    \vspace{2pt}
    {\footnotesize 作成日: 2026年10月28日 \quad 提出者: IbukiSakamoto}\par
    \vspace{2pt}
    \hrule
  \end{center}
  \vspace{-10pt}
]

\section{開発の概要と目的}
本開発では，Java 21，Spring Boot 3，およびMySQLを用いて学習用ECサイトの最小構成Webアプリケーションを構築した．MVCアーキテクチャ，Spring Data JPAによるO/Rマッピング，セッション管理，トランザクション制御といったWeb開発の中核技術の実践的な理解を目的としている．

\section{使用技術とシステム構成}
\begin{itemize}\setlength{\itemsep}{-0.5pt}\setlength{\parskip}{0pt}
  \item \textbf{言語・FW}: Java 21 (OpenJDK), Spring Boot 3.3.4
  \item \textbf{Web・DB}: Spring MVC, Thymeleaf, MySQL 8.4
  \item \textbf{データアクセス}: Spring Data JPA (Hibernate)
  \item \textbf{ビルド}: Apache Maven 3.9
  \item \textbf{通信構成}: ホストPC(Windows)からの直接アクセスを可能とするため，コンテナの80番ポートへの通信を\texttt{socat}によりSpring Boot(8080番)へ転送する構成を採用した．
\end{itemize}

\section{実装した主な機能}
画面表示はすべて日本語に統一し，以下の機能を実装した．
\begin{itemize}\setlength{\itemsep}{-0.5pt}\setlength{\parskip}{0pt}
  \item \textbf{商品一覧}: 全商品のカード表示および在庫状態バッジ表示（在庫あり／残りわずか／在庫切れ）．
  \item \textbf{商品検索・ソート}: 商品名・説明文の部分一致キーワード検索，価格の昇順・降順ソート．
  \item \textbf{カート}: セッション単位での商品保持，数量変更，削除，カート全消去，合計金額算出．
  \item \textbf{初期データ投入}: アプリケーション起動時に\texttt{CommandLineRunner}により商品8件を自動登録．
\end{itemize}

\section{注文処理と在庫管理の仕組み}
注文処理はサービス層で\texttt{@Transactional}を付与して実行し，整合性を担保している．
カート内各商品の最新在庫数をDBから検証し，不足時は\texttt{OutOfStockException}を送出してロールバックする．
在庫が確保できる場合，在庫数を減算更新した上で\texttt{Order}および\texttt{OrderItem}（購入時点の単価・数量・小計）を永続化し，カートを初期化して注文完了画面を表示する．

\section{Antigravity CLIを活用した開発手順}
開発はGoogle Antigravity CLIを用い，エージェントとの対話型指示により段階的に進めた．
\begin{enumerate}\setlength{\itemsep}{-0.5pt}\setlength{\parskip}{0pt}
  \item 環境調査（Ubuntu環境，Java・Maven・MySQL導入確認）
  \item OpenJDK 21，Maven，MySQLのインストールとDB初期化
  \item Spring Bootプロジェクトの骨格・各層コード作成
  \item ビルドおよび動作検証と不具合修正
  \item \texttt{socat}による80番ポート転送の設定
  \item 機密情報の排除確認，\texttt{.gitignore}整備，Gitコミット
  \item GitHub CLIを用いた安全なデバイス認証とGitHub公開
\end{enumerate}

\section{動作確認の結果}
Windows環境のブラウザから\texttt{http://localhost/}にアクセスし，以下の正常動作を確認した．
\begin{itemize}\setlength{\itemsep}{-0.5pt}\setlength{\parskip}{0pt}
  \item 商品一覧の表示，キーワード検索および価格順ソート
  \item カートへの追加，数量更新，小計・合計の正確な計算
  \item 注文者名入力から注文確定，注文完了画面の表示
  \item 注文確定に伴う在庫数の正確な減算，および在庫不足時におけるカート追加・注文拒否のバリデーション
\end{itemize}

\section{今後の改善点}
学習用の最小構成として開発したため，実運用に向けて以下の機能拡張が挙げられる．
\begin{itemize}\setlength{\itemsep}{-0.5pt}\setlength{\parskip}{0pt}
  \item Spring Securityを導入した会員登録・認証機能
  \item 外部決済サービスとの連携および注文履歴一覧
  \item 管理者向けの商品登録・編集・在庫管理画面
\end{itemize}

\section{GitHubリポジトリURL}
ソースコードおよびREADMEは以下で公開されている．\\
{\footnotesize \url{https://github.com/IbukiSakamoto/spring-boot-ecsite}}

\end{document}
